% БҮЛЭГ 1 — АСУУДЛЫН ТОМЬЁОЛОЛ (УДИРТГАЛ)

\chapter{Удиртгал}
\label{Chapter1}

Энэхүү бүлэгт дипломын судалгааны үндэслэл, шийдэж буй асуудлын тодорхой
томьёолол, SMART зарчмын дагуу томьёолсон зорилго ба зорилтууд, ажлын
таамаглал, хамрах хүрээ болон шинжлэх ухаан, нийгэмд оруулах хувь нэмрийг
нарийвчлан танилцуулна.

%-------------------------------------------------------------------------------
\section{Судалгааны үндэслэл}

Интернэтийн хэрэглэгчдийн тоо дэлхий даяар урьд өмнө байгаагүй хурдацтай
нэмэгдэж байгаа бөгөөд цахим мэдээллийн орчин нийгмийн харилцааны хэлбэрийг
бүрэн өөрчилсөн. DataReportal-ийн (2024) мэдээллийн дагуу \cite{datareportal2024}
Монгол Улсын интернэт хэрэглэгчдийн тоо нийт хүн амын 80 гаруй хувийг эзэлж байгаа бөгөөд
Facebook, X (Twitter), TikTok зэрэг платформуудын хэрэглэгчид жил бүр хурдацтай
өсч байна. Цахим орчинд тархаж буй мэдээлэл, харилцааны хэмжээ нэмэгдэхийн
хэрээр хэрэглэгчид, ялангуяа өсвөр насныхан, эмзэг, хортой агуулгад
өртөх эрсдэл ч нэгэн зэрэг нэмэгдэж байна.

Монголын цахим орчинд --- тухайлбал news.mn зэрэг мэдээний сайтын сэтгэгдлийн
хэсэг, Facebook-ийн нийтийн хуудас болон бүлгүүд, X (Twitter)-ийн нийтийн
нийтлэлүүдэд --- кибер дарамт, сөрөг агуулга, хорлолтой дезинформаци, хэт
аггрессив хандлага зэрэг олон хэлбэрийн сөрөг агуулга тархаж байна. Олон
улсын судалгаануудад \cite{kowalski2014,vogels2022} кибер дарамт нь өсвөр насны
хүүхдийн сэтгэл зүйн тогтвортой байдал, нийгмийн харилцаа, академик
гүйцэтгэлд ноцтой сөргөөр нөлөөлдгийг нотолсон. Монгол Улсад энэ асуудалд
хариулах техникийн шийдэл, тухайлбал автомат контент шүүлтийн систем,
хөгжлийн түвшин хангалтгүй байгааг тэмдэглэх нь зүйтэй.

Гэхдээ аливаа сөрөг агуулгыг цогцоор нь хориглох нь нийгмийн хувьд буруу
шийдэл болно. Бүтээлч шүүмжлэл, асуудалд суурилсан шинжилгээ, бодлогыг
эсэргүүцэх тайлбар зэрэг нь ардчилсан нийгэм, иргэний соёлын чухал суурь
болдог. ``Positive-only'' орчин нь нийгмийн бодит байдлаас тасарсан хиймэл
дискурсийг бий болгох аюултай бөгөөд үзэл бодлоо чөлөөтэй илэрхийлэх нь
үндсэн хүний эрхэд хамаарна. Тиймд автомат шүүлтийн систем нь зөвхөн бүх
сөрөг агуулгыг устгах зорилготой байж болохгүй, харин хоёр эсрэг туйлын
агуулгыг --- нэг нь хориглох, нөгөөг нь хамгаалах --- нарийн ялгаж чаддаг
байх шаардлагатай.

%-------------------------------------------------------------------------------
\section{Асуудлын томьёолол}

Монгол хэлний тусгайлсан NLP судалгааны хүрээнд Хортой сөрөг болон Бүтээлч шүүмжлэл
агуулгыг ялгах систем одоогоор боловсруулагдаагүй байгаа нь судалгааны
орхигдсон чиглэлийг бүрдүүлж байна. Өмнөх ажлуудын дийлэнх нь гурван ангилалт
(Эерэг / Сөрөг / Саармаг) сентимент шинжилгээгээр хязгаарлагдаж,
``Сөрөг'' ангилалын дотоодод байх хоёр эсрэг шинж чанарыг --- устгах
шаардлагатай хортой агуулга ба хамгаалах шаардлагатай бүтээлч шүүмжлэл ---
ялган тодорхойлоогүй байна. Энэхүү судалгааны ажил тэрхүү орхигдсон чиглэлийг нөхөх
зорилготой бөгөөд нийгмийн болон техникийн ач холбогдлоороо тодорхой хувь
нэмэр оруулна.

%-------------------------------------------------------------------------------
\section{Судалгааны зорилго}

Монгол хэл дээрх цахим орчны агуулгыг олон ангилалтаар шүүх, тухайлбал Хортой сөрөг
болон Бүтээлч шүүмжлэлийг нарийн ялгаж чаддаг хиймэл оюунд суурилсан
NLP системийг боловсруулах нь энэхүү судалгааны гол зорилго юм. Тус систем нь
Монгол хэлний agglutinative морфологийн онцлогт тохирсон өгөгдлийн
боловсруулалтын дамжлага, MN-BERT загварын нарийн тохируулга (fine-tuning),
болон Gradio-д суурилсан хэрэглэгчийн интерфейсийг нэгтгэнэ.

Зорилгыг SMART зарчмын дагуу томьёолбол:

\begin{itemize}
  \item \textbf{Specific (тодорхой):} 4 ангилалтай (Эерэг, Саармаг, Хортой сөрөг,
    Бүтээлч шүүмжлэл) Монгол кирилл текстийн ангилагч
    системийг хоёр шатлалт архитектуртай зохион бүтээх.
  \item \textbf{Measurable (хэмжигдэхүйц):} End-to-end macro F1 оноо $\geq 0.80$;
    нарийвчлал зорилтот түвшин $\geq 90\%$ (Ш-1, 4 оноо), үнэлгээний рубрикийн
    дагуу $80$--$89\%$ нь 3 оноогоор тооцогдоно.
  \item \textbf{Achievable (биелэгдэхүйц):} MN-BERT-ийн pre-trained жинг
    суурь болгон нарийн тохируулга хийж, ангиллын тэнцвэргүй байдлыг суурь
    загваруудад SMOTE-оор, эцсийн нэг шатлалт MN-BERT загварт ангиллын жин
    оногдуулсан Focal Loss ($\gamma{=}2.0$) болон WeightedRandomSampler-аар
    шийдвэрлэх.
  \item \textbf{Relevant (хамааралтай):} Монгол хэлний NLP-ийн одоогийн
    ``Эерэг/Саармаг/Сөрөг'' гурван ангилалт суурь дээр Хортой сөрөг vs
    Бүтээлч шүүмжлэл ялгалт нэмэлтээр оруулах орхигдсон чиглэлтэй шууд холбоотой.
  \item \textbf{Time-bound (хугацаатай):} 2025 оны 9 сараас 2026 оны 5 сар
    хүртэлх дипломын ажлын хугацаанд хэрэгжүүлэх.
\end{itemize}

%-------------------------------------------------------------------------------
\section{Судалгааны зорилтууд}

Дипломын ажлын хүрээнд дараах зорилтуудыг дараалсан байдлаар хэрэгжүүлнэ.

\begin{enumerate}
  \item Найман нийтэд нээлттэй эх сурвалжаас (news.mn, gogo.mn, IKON.mn,
        E-Mongolia, Facebook-ийн нээлттэй групп зэрэг; бүлэг~\ref{Chapter4},
        Хүснэгт~\ref{tab:data_full_corpus}) Монгол хэл дээрх сэтгэгдэл болон
        нийтлэлүүдийг цуглуулах, цэвэрлэх, стандартжуулах ажлыг гүйцэтгэнэ.

  \item Цуглуулсан өгөгдлийг Эерэг, Саармаг, Хортой сөрөг, Бүтээлч шүүмжлэл
        гэсэн дөрвөн ангилалд гар аргаар шошголох (manual
        annotation) бөгөөд гурван шошгологчийн хоорондын нийцлийг Fleiss' $\kappa$
        коэффициентээр хэмжинэ.

  \item Монгол хэлний agglutinative онцлогт тохирсон урьдчилсан боловсруулалтын
        (preprocessing) pipeline --- кирилл шүүлт, нормчлол, токенчлол (стоп үг
        хасах ба stemming нь зөвхөн TF-IDF суурь загварт) --- боловсруулна.

  \item Нэг шатлалт (flat) ба хоёр шатлалт (two-stage) ангиллын архитектурыг
        зэрэг зохион байгуулж харьцуулна (хоёр шатлалт: Stage~1 ---
        Эерэг/Саармаг/Сөрөг; Stage~2 --- Сөрөг агуулгыг Хортой сөрөг ба
        Бүтээлч шүүмжлэл болгон нарийвчлах), туршилтын үр дүнд тулгуурлан
        эцсийн архитектурыг сонгоно.

  \item MN-BERT загварыг нарийн тохируулга хийж Logistic Regression, Naive Bayes,
        SVM зэрэг суурь загваруудтай харьцуулан статистик аргаар үнэлнэ.

  \item Accuracy, Precision, Recall, $F_1$-score, Confusion Matrix зэрэг
        үнэлгээний үзүүлэлтүүдэд тулгуурлан загваруудын гүйцэтгэлийг
        тодорхойлно.

  \item Gradio framework ашиглан хэрэглэгчийн интерфейстэй прототип
        системийг хэрэгжүүлж, туршилтын орчинд үнэлнэ.
\end{enumerate}

%-------------------------------------------------------------------------------
\section{Таамаглал}

Судалгааны гол таамаглал нь: \textit{``Монгол хэлний agglutinative онцлогт
тохирсон урьдчилсан боловсруулалтын дамжлага дээр суурилан, MN-BERT pre-trained загварыг
хоёр шатлалт (two-stage) архитектурт нарийн тохируулга хийж ашиглавал, TF-IDF
суурьтай уламжлалт машин сургалтын загваруудаас (LR, NB, SVM)
дээгүүр macro F1 гаргаж, Хортой сөрөг болон Бүтээлч шүүмжлэлийг
найдвартай ялгах боломжтой.''}

Дэд таамаглалууд:
\begin{itemize}
  \item H$_1$: MN-BERT нарийн тохируулгатай загвар нь TF-IDF + LR/NB/SVM-ээс macro F1 оноо
    хэмжүүрээр дээгүүр гүйцэтгэл үзүүлнэ.
  \item H$_2$: Ангиллын тэнцвэржүүлэлт (SMOTE суурь загварт; Focal Loss + жин
    $\alpha$ + WeightedRandomSampler эцсийн нэг шатлалт MN-BERT-д) ба өгөгдлийн
    чанаржуулалт нь нийт гүйцэтгэлийг $\geq 10\%$ сайжруулна.
  \item H$_3$: Хоёр шатлалт архитектур нь Хортой сөрөг vs Бүтээлч шүүмжлэлийг flat
    4-ангиллын загвараас илүү нарийвчилсан тооцоолох pipeline өгнө (хэдий
    каскад алдааны эрсдэлтэй ч).
\end{itemize}

%-------------------------------------------------------------------------------
\section{Судалгааны хамрах хүрээ ба хязгаарлалт}

Энэхүү судалгаа нь дараах хэмжээст орон зайд явагдана. Хэлний хувьд зөвхөн
Монгол кирилл бичгийн текст хамрагдана. Латин үсгийн бичвэр, монгол-англи
холимог (code-switched) текст, латин үсгийн хэлбэртэй Монгол бичлэг (``bn'', ``yu ve'',
``zgr'' зэрэг), болон бусад кирилл бус бичиглэлийн хэлбэрүүд датасетаас
тусгайлан хасагдана. Энэхүү шийдэл нь MN-BERT-ийн кирилл корпус дээр суурилсан токенчлолын
тогтвортой байдлыг хамгаалахын зэрэгцээ шошиглолтын нэгдсэн стандартыг
баримтлахад зайлшгүй шаардлагатай.

Платформын хувьд мэдээний сайтуудын (news.mn, gogo.mn, IKON.mn гэх мэт) нийтийн
сэтгэгдлийн хэсэг болон Facebook-ийн нээлттэй группийн сэтгэгдэл --- нийт найман
эх сурвалж (бүлэг~\ref{Chapter4}, Хүснэгт~\ref{tab:data_full_corpus}) ---
хамрагдана. Хугацааны хувьд 2025--2026 оны эхний нийтлэлүүдийг
ашиглана. Контентийн хувьд зөвхөн текстэн өгөгдлийг анализад оруулах бөгөөд
зураг, видео, эможи зэргийг анхдагч загварт оруулахгүй --- эдгээрийг ирээдүйн
судалгааны чиглэл болгон тодорхойлно. Системийн хувьд прототип загвар
боловсруулахаар тооцох бөгөөд үйлдвэрлэлийн орчинд байршуулах нь энэхүү ажлын
хүрээнд багтахгүй.

%-------------------------------------------------------------------------------
\section{Ажлын шинэлэг тал ба ач холбогдол}

Дипломын ажлын шинэлэг тал нь, бидний мэдэхээр (бүлэг~\ref{Chapter2}-ийн
холбогдох судалгааны шинжилгээнд тулгуурлан), Монгол хэлний цахим орчинд
Хортой сөрөг ба Бүтээлч шүүмжлэлийг тусгайлан ялгах \textit{анхны} системчилсэн
оролдлого, мөн нэг шатлалт ба хоёр шатлалт архитектурын харьцуулсан судалгаа
явуулж буйд оршино. Үүнээс гадна дараах гурав нь салбарын одоогийн төлөв байдалд тодорхой хувь нэмэр болно:
\begin{enumerate}
  \item Монгол хэлний agglutinative онцлогт тохирсон нэгдсэн урьдчилсан боловсруулалтын дамжлага,
  \item MN-BERT-ийн нарийн тохируулгыг FocalLoss болон cosine scheduler-тэй хослуулсан туршилт,
  \item аннотаторуудын хоорондын нийцлийн хэмжилтийг (Fleiss' $\kappa = 0.72$) ил тодоор тайлагнасан, дөрвөн ангилалт шошгологдсон өгөгдлийн багц.
\end{enumerate}

Нийгмийн ач холбогдлын хувьд уг систем нь Монгол цахим орон зайд хортой агуулга,
кибер дарамт зэрэг сөрөг агуулгыг хянан илрүүлэх хэрэгсэл болохын зэрэгцээ
иргэдийн хуулийн хүрээнд хамгаалагдсан бүтээлч шүүмжлэлийн эрхийг хамгаалах
ёс зүйн хэв загвар бүрдүүлнэ.

%-------------------------------------------------------------------------------
\section{Бүлгийн дүгнэлт}

Монгол хэлний цахим орчинд хортой болон бүтээлч агуулгыг автоматаар ялган
шүүх асуудал нь зөвхөн техникийн биш, нийгэм, иргэний болон хүний эрхийн
чухал ач холбогдолтой. Өмнөх судалгааны ажлуудад энэ ялгалт хийгдэж
байгаагүй нь тодорхой судлагдаагүй асуудлыг бий болгосон. Энэхүү бүлэгт
асуудлын томьёолол, SMART зорилго, гурван дэд таамаглал, хамрах хүрээ,
шинэлэг тал болон нийгмийн ач холбогдлыг тодорхой тогтоов. Дараагийн
бүлэгт (бүлэг~\ref{Chapter2}) NLP, мэдрэмжийн шинжилгээ, BERT архитектур
болон холбогдох ажлуудын онолын суурийг нарийвчлан авч үзнэ.
